\chapter{Scope, research questions and project history}
\label{ch:scope}
\section{Original objective and changing emphasis}
The project began as a comparison of compact contrastive and generative vision--language models for public-space anomaly recognition. The requested targets are smoking, littering, loitering, rough sleeping and unattended personal items. The requested model families are \modelm, \capstonesig, SmolVLM-500M-Instruct, MiniCPM-V 4.0 and VadCLIP. The early repository emphasized a constrained CPU deployment setting. The subsequent user direction explicitly prioritized model quality, allowed available GPU compute without pacing, and eventually removed the need for further latency testing.

That change affects interpretation of the work. Accuracy experiments should use available compute to test the models properly; a slow CPU replay does not make the quality test more valid. Conversely, CPU latency requires a separately controlled protocol. The later experiments therefore use GPU training and extraction where applicable, with no artificial camera-rate sleeps. The historical CPU pilot remains part of the record, but no new latency study is added after the change in priorities.

The research questions become progressively more specific. First, can the frozen models produce useful categorical decisions using common prompts and images? Second, does a larger and more carefully reviewed corpus change those conclusions? Third, can supervised binary heads and actual visual adaptation improve task discrimination? Fourth, do crops, explicit temporal order, geometry and controlled weight blending transfer better to other data sources? Fifth, can causal item history and event accounting address the gap between high image AP and excessive alerts? The 17 October experiments then isolate ensembling, source-wise operating points, stationarity smoothing, alarm memory and image tiling. The 19 October study asks whether language-layer LoRA can improve a task-conditioned SmolVLM decision function across all five definitions, while preserving a matched original-model control and independently repeating the winning learning rate.

\section{What each target requires}
\begingroup\footnotesize
\setlength{\LTpre}{5pt}\setlength{\LTpost}{8pt}
\begin{longtable}{@{}>{\raggedright\arraybackslash}p{0.1680\linewidth}>{\raggedright\arraybackslash}p{0.6720\linewidth}@{}}
\caption[Definitions]{Recorded visible-evidence target definitions.}\label{tab:definitions}\\
\toprule
\textbf{Label} & \textbf{Definition} \\
\midrule\endfirsthead
\multicolumn{2}{l}{\textit{Table \thetable\ continued}}\\
\toprule
\textbf{Label} & \textbf{Definition} \\
\midrule\endhead
\midrule\multicolumn{2}{r}{\textit{Continued on next page}}\\\endfoot
\bottomrule\endlastfoot
normal & A public area with ordinary activity and none of the defined anomalies visibly present. \\
littering & A person visibly discarding rubbish on the ground instead of using a bin. \\
loitering & A person remaining in the same public area for an extended time without an apparent activity. Duration requires evidence across time. \\
rough sleeping & A person apparently sleeping or lying down with bedding in a public area. Describe visible evidence only; do not infer housing status. \\
smoking & A person visibly smoking a cigarette, cigar, or pipe, with a smoking object near the mouth or visible smoke. \\
unattended items & A bag, box, or other personal item apparently left behind without its owner nearby. Ownership and duration require evidence across time. \\
\end{longtable}\endgroup

The definitions are the repository definitions accepted during the work. There is no newly imposed universal loitering duration. Source-provided event criteria remain explicit. In particular, AVSS's 30-second stationary-object criterion is a source-specific diagnostic, not a retrospective relabelling of every historical unattended interval.

Several targets require evidence absent from a single global image. A person near a cigarette-sized object is spatially subtle; a person remaining in an area is a temporal identity problem; a bag becoming unattended is a relationship and duration problem. Apparent public sleep needs visible posture/context and, for a persistent event, time evidence. The model interface should therefore permit unknown when duration, ownership or visibility is insufficient. A scene classifier alone cannot establish object localization or ownership.

\section{Model families and comparison boundaries}
\begingroup\footnotesize
\setlength{\LTpre}{5pt}\setlength{\LTpost}{8pt}
\begin{longtable}{@{}>{\raggedright\arraybackslash}p{0.1515\linewidth}>{\raggedright\arraybackslash}p{0.1377\linewidth}>{\raggedright\arraybackslash}p{0.2754\linewidth}>{\raggedright\arraybackslash}p{0.2754\linewidth}@{}}
\caption[Models]{Models and comparison scope; pinned revisions appear in the appendix.}\label{tab:models}\\
\toprule
\textbf{Model} & \textbf{Family} & \textbf{Checkpoint / feature source} & \textbf{Role} \\
\midrule\endfirsthead
\multicolumn{4}{l}{\textit{Table \thetable\ continued}}\\
\toprule
\textbf{Model} & \textbf{Family} & \textbf{Checkpoint / feature source} & \textbf{Role} \\
\midrule\endhead
\midrule\multicolumn{4}{r}{\textit{Continued on next page}}\\\endfoot
\bottomrule\endlastfoot
Mobile\allowbreak{}CLIP & openclip & \path{timm/MobileCLIP2-S2-OpenCLIP} & Frozen image comparison; CLIP later adapted \\
SigLIP & siglip & \path{google/siglip2-base-patch16-224} & Frozen image comparison; CLIP later adapted \\
SmolVLM & vlm & \path{HuggingFaceTB/SmolVLM-500M-Instruct} & Original prompted comparison; 19 Oct language-layer LoRA with binary scoring \\
MiniCPM & minicpm & \path{openbmb/MiniCPM-V-4} & Prompted image comparison only \\
VadCLIP & Supervised video adapter & Frozen CLIP-B/16 features & Separate offline smoking protocol \\
\end{longtable}\endgroup

The CLIP-family adapters produce continuous text-image scores. For the original shared image protocol, each model compares the same selected inputs and definitions. Generative VLMs answer a deterministic classification prompt and may emit a valid label, explicit unknown or invalid response. Hard labels do not provide a continuous ranking score from which a meaningful AP or AUROC can be invented.

VadCLIP is a separate supervised offline smoking adaptation over frozen CLIP-B/16 features. It receives sequence supervision, uses a different split and can use context within an offline sequence. It is therefore informative about a temporal supervised baseline, but it is not an equal-supervision five-task competitor in the image-model table. The 15--17 October learning studies adapt the two CLIP backbones and fit subsequent heads. On 19 October, a separate study actually adapts SmolVLM language projections with LoRA under matched binary prompting. MiniCPM and the historical VadCLIP adapter receive no additional training in that round.

\section{The early small benchmark and why it was superseded}
The historical 11 October narrative describes an 80-image benchmark with 16 examples each of littering, loitering, smoking, rough sleeping and normal, and no unattended-item examples. It reported a six-label multiclass comparison and various CPU and VadCLIP ablations. Later audits established that some apparent target labels were only person, action or synthetic posture proxies. The historical binary anomaly prevalence was 64/80; an always-anomaly classifier already obtains binary F1 of approximately 0.889. Such a base rate can make a high F1 look more convincing than the event evidence supports.

The report is retained as historical context, not used as the main five-event conclusion. The subsequent validation summary explicitly narrows its scope to a real MobileCLIP CPU pipeline check, rather than a full-model/full-event ranking. This is an important correction in the project: unsupported event labels and incomplete model execution must not be carried forward merely because an earlier report presented a leaderboard.

The 12 October run originally open in the IDE is another distinct artifact. It is marked execution status \texttt{error}; its prepared smoking cohort has only 236 frames, and the MiniCPM smoking result is missing. Some saved response counts include three repetitions of the same selected input. A repeated inference is not a new labelled sample. This run is useful for diagnosing caps, CPU execution and response behavior, but the completed uncapped GPU study replaces it as the principal baseline comparison.

\subsection{Correctness problems identified before the new quality claims}
The archived 11 October audit records 15 findings, not merely a request to add more samples. It reproduces fractional-positive labels being truncated to zero after temporal compression: a mixed segment label .5 becomes integer zero at evaluation. Changing visual length therefore changes both targets and prevalence, invalidating a clean length-ablation interpretation. It also reproduces a precision--recall threshold-index mismatch, so the reported threshold can disagree with the reported F1; subsequent evaluation recomputes decisions at actual development-selected cutoffs.

A CPU ``best checkpoint'' created using a detached CPU tensor without a clone shares storage with ongoing training weights. Later optimizer steps can change the saved in-memory best state. The new training paths retain cloned snapshots and checkpoint identity checks. Other audit concerns include validation reused for epoch/branch/threshold selection without a final test, padded frames influencing real predictions, cache reuse without complete image/annotation/weight identity, and accepting stale notebook summaries after incomplete configuration checks.

The image-label audit also shows that normal meant not-smoking in one corpus, not verified absence of every event, while proxy class membership could be trivially confounded with the source. A post-hoc metadata-only dimension rule is reported as reaching 80\% accuracy on the old 80-image bundle; this is a source-confounding demonstration, not a legitimate held-out detector baseline. The historical macro-F1 class denominator could change with predicted classes, so the old claim of a fixed five-sixths upper bound was not a reliable scoring description. Abstentions were excluded from some conditional scores, and row-wise rather than source-group bootstrap exaggerated independence.

These findings explain the later emphasis on source-qualified binary tasks, fixed supported label sets, all-input coverage, cloned checkpoints, raw-frame mappings, development-frozen cutoffs, immutable experiment directories and independent metric replay. The appendix preserves the original finding inventory and the historical VadCLIP ablation records as superseded evidence. Their displayed APs are not promoted into the corrected final leaderboard.
\evidence{reports/anomaly_benchmark_report.md}
\evidence{reports/archive/2026-10-11/evaluation_validation_summary.json}
\evidence{reports/archive/2026-10-11/evaluation_audit.md}
\evidence{results/evaluation_20261012T165345_130231Z/report.md}

\section{Experiment chronology}
\begingroup\footnotesize
\setlength{\LTpre}{5pt}\setlength{\LTpost}{8pt}
\begin{longtable}{@{}>{\raggedright\arraybackslash}p{0.1273\linewidth}>{\raggedright\arraybackslash}p{0.2036\linewidth}>{\raggedright\arraybackslash}p{0.5091\linewidth}@{}}
\caption[Chronology]{Distinct stages of the recorded project.}\label{tab:chronology}\\
\toprule
\textbf{Dates} & \textbf{Stage} & \textbf{Completed evidence / important qualification} \\
\midrule\endfirsthead
\multicolumn{3}{l}{\textit{Table \thetable\ continued}}\\
\toprule
\textbf{Dates} & \textbf{Stage} & \textbf{Completed evidence / important qualification} \\
\midrule\endhead
\midrule\multicolumn{3}{r}{\textit{Continued on next page}}\\\endfoot
\bottomrule\endlastfoot
11 Oct & Early benchmark and audit & 80 images; proxy labels and no unattended class. Historical, superseded. \\
12 Oct & Corrected CPU-era run & 236 smoking frames; run status error; MiniCPM smoking missing. \\
13 Oct & Full GPU comparison & 13,400 predictions across four models; separate supervised VadCLIP and CPU pilot. \\
14 Oct & Expanded accuracy & 55,156 predictions; 1,770 unattended endpoints and qualified weak diagnostics. \\
15--16 Oct & Visual fine-tuning & Two backbones x three seeds; six actual changed-visual-weight runs. \\
16 Oct & Strategy screen & 330 configurations, new-source transfer and provided thermal trajectories. \\
16 Oct & Causal video pilot & 50 heads; real tracks/regions; two clean new AVSS events, one censored. \\
17 Oct & Focused improvements & 38 policies; state/memory ablations; 830 tiled-detector GPU views. \\
18 Oct & First full technical report & 83-page account through 17 October; no new training for documentation. \\
19 Oct & SmolVLM LoRA and updated documentation & Three trained adapters; ten held-out cohorts; repeat-seed audit; no latency testing. \\
\end{longtable}\endgroup

The chronology should be read as a sequence of distinct studies. The earlier test images become historically inspected after their first evaluation. Later development-only selection is still preferable to choosing on test results, but repeated experimentation against familiar source families cannot restore an untouched test claim. The report consequently uses ``retrospective'' whenever later changes reuse those results.

\section{What was preserved}
Every completed study has its own result directory, selection records and raw predictions. Original AMP outputs remain beside corrected float32 evaluations. Earlier source-prefix video fits remain in a diagnostic subdirectory after the known-film guard. The initial incorrect tiled-metric export remains explicitly labelled invalid after an aggregation bug was corrected from the identical cached detections. Earlier slide decks remain historical; the 17 October deck consolidates 32 slides without replacing their source results. The 19 October edition extends the technical paper and presentation to include the completed SmolVLM study while retaining the historical versions. This preservation makes the progression auditable and prevents a later narrative from erasing regressions.
